\clearpage\begingroup
\SourceLeaf{290}\NoSourceHead
\begin{center}\Large ÜBER DAS KREISMIKROMETER.\par\bigskip\large\textls[120]{BRIEFWECHSEL.}\end{center}
\begin{center}\small
\rule{\textwidth}{.35pt}\\[.2em]
[\textsc{Wilhelm Olbers}, Sein Leben und seine Werke, II\textsubscript{1}, Berlin 1900, S. 402. 409 und II\textsubscript{2}, Berlin 1909, S. 69.]\\[-.2em]
\rule{\textwidth}{.35pt}
\end{center}
\begin{center}\textsc{Gauss} an \textsc{Olbers}, Göttingen, 5. Januar 1808.\end{center}
\noindent\mbox{. . . . . . . . . . . . . . .}\par
\noindent\hspace*{\parindent}Mein Verfahren zur Reduction der Kreismikr[ometer-]Beobacht[ungen] mit Rücksicht auf Refraction gründet sich darauf, dass in dem kreisförmigen Gesichtsfelde sich vom Himmel ein Ausschnitt zeigt, der als elliptisch angesehen werden kann: die grosse Axe dieser Ellipse liegt in dem Vertikalkreise; diese Voraussetzung weicht nur ganz unmerklich von der Wahrheit ab. Es sei
\begin{center}\begin{tabular}{rl}
$2r$&wahrer Halbmesser des Gesichtsfeldes,\\
$2kr$&grosse (vertikale) Axe der Ellipse,\\
$2hr$&kleine (horizontale) Axe der Ellipse,
\end{tabular}\end{center}
\noindent so werden $h$, $k$ von der Höhe abhängige Faktoren, etwas weniges grösser als $1$ sein. Setzt man für den sich im Mittelpunkt des Fernrohrs zeigenden Stern die wahre Zenithdistanz $=z$, die scheinbare $=z'=z-\rho$, so ist (hinreichend genau):
\[
k=\frac{dz}{dz'}=\frac{1}{1-\dfrac{d\rho}{dz}},\qquad h=\frac{\sin z}{\sin z'}\mathpunct{\cdot}
\]
\noindent\hspace*{\parindent}Setzt man nun $\surd\bigl(1-\frac{hh}{kk}\bigr)=\lambda$ (wo $\lambda$ die Excentricität der Ellipse bedeuten wird), so hat man folgende Formeln:
\begin{center}\begingroup\small\setlength{\tabcolsep}{3pt}
\begin{tabular}{@{}lr@{}}
Zeit des Eintritts und Austritts (in Secunden)\quad\mbox{. . . . . . . . . . . . .}&$t,t'$\\
\multicolumn{2}{@{}l@{}}{Zeit, wo dem $\ast$ und dem Mittelpunkte des Fernrohrs einerlei Stunden-}\\
\quad winkel entspricht\quad\mbox{. . . . . . . . . . . . . . . . . . . . . . . . . .}&$\Theta$\\
Declination des Mittelp[un]kts\quad\mbox{. . . . . . . . . . . . . . . . . . . . . .}&$D$\\
\quad »\quad »\quad Sterns\,\mbox{. . . . . . . . . . . . . . . . . . . . . . . . . . . . . .}&$\delta$\\
Positionswinkel (positiv westlich)\,\mbox{. . . . . . . . . . . . . . . . . . . . .}&$p$
\end{tabular}\endgroup\end{center}
\[\frac{2rh}{15}=R.\]
\clearpage\SourceLeaf{291}\HistoricalHead{286}{L}{PRAKTISCHE UND SPHÄRISCHE ASTRONOMIE. BRIEFWECHSEL.}
\noindent\hspace*{\parindent}Man setze
\[\mathrm{I.}\quad\frac{(t'-t)\cos\delta}{R}\surd(1-\lambda\lambda\sin p^2)=\cos\psi\]
\noindent so ist
\begingroup\small
\[\mathrm{II.}\quad\delta=D\pm kr\sin\psi\surd(1-\lambda\lambda\sin p^2)=D\pm\frac{k}{h}\cdot hr\sin\psi\surd(1-\lambda\lambda\sin p^2)\]
\[\mathrm{III.}\quad\Theta=\tfrac12(t+t')+\frac{\delta-D}{15\cos\delta}\cdot\frac{\lambda\lambda\sin p\cos p}{1-\lambda\lambda\sin p^2}\mathpunct{\cdot}\]
\endgroup
\noindent $h$ kann man für alle Höhen als constant betrachten; ich finde
\[\begin{array}{rclll}
\log h&=&0{,}00012&\text{von }z=0^0&\text{bis }82^0\\
\log h&=&0{,}00011&\text{von }z=83^0&\text{bis }85^0\\
\log h&=&0{,}00010&\text{für }z=86^0&\\
\log h&=&0{,}00009&\text{für }z=87^0.&
\end{array}\]
\noindent Es ist also nur eine Tafel für $\frac{k}{h}$ und $\lambda$ erforderlich; hier eine solche, flüchtig berechnet:
\begin{center}\begingroup\fontsize{9.5}{12}\selectfont\setlength{\tabcolsep}{3pt}\renewcommand{\arraystretch}{1.13}
\begin{tabular}{r|r|r||r|r|r||r|r|r||r|r|r}
\hline\hline
$z$&$\log\frac{k}{h}$&$\log\lambda$&$z$&$\log\frac{k}{h}$&$\log\lambda$&$z$&$\log\frac{k}{h}$&$\log\lambda$&$z$&$\log\frac{k}{h}$&$\log\lambda$\\[.2em]
\hline\hline
$0^0$&$0{,}00000$&$-\infty$&$22^0$&$0{,}00002$&$7{,}982$&$44^0$&$0{,}00011$&$8{,}360$&$66^0$&$0{,}00061$&$8{,}723$\\
1&&$6{,}618$&23&2&$8{,}003$&45&12&375&67&67&743\\
2&&$6{,}919$&24&2&024&46&13&390&68&73&764\\
3&&$7{,}095$&25&3&044&47&14&405&69&81&786\\
4&&220&26&3&064&48&15&420&70&$0{,}00090$&808\\
5&&317&27&3&083&49&16&435&71&$0{,}00100$&832\\
6&&397&28&3&101&50&17&451&72&112&856\\
7&&463&29&4&119&51&19&466&73&126&882\\
8&&523&30&4&137&52&20&481&74&143&908\\
9&&575&31&4&154&53&21&497&75&162&936\\
10&&622&32&5&171&54&23&513&76&186&966\\
11&$0{,}00000$&664&33&5&188&55&25&529&77&215&$8{,}997$\\
12&$0{,}00001$&703&34&6&204&56&27&545&78&251&$9{,}030$\\
13&&739&35&6&220&57&29&561&79&295&065\\
14&&772&36&6&236&58&31&578&80&353&104\\
15&&804&37&7&252&59&34&595&81&425&144\\
16&&833&38&7&268&60&36&612&82&519&187\\
17&&861&39&8&283&61&39&629&83&651&235\\
18&&887&40&9&299&62&43&647&84&$0{,}00836$&288\\
19&$0{,}00001$&912&41&$0{,}00009$&314&63&47&665&85&$0{,}01081$&$9{,}343$\\
20&$0{,}00002$&937&42&$0{,}00010$&329&64&51&684&&&\\
21&$0{,}00002$&$7{,}960$&43&$0{,}00011$&$8{,}345$&65&$0{,}00055$&$8{,}703$&&&
\end{tabular}\endgroup\end{center}
\clearpage\endgroup
